\clearpage
\begin{explanationbox}
\paragraph{同样在原点不可导：$x^{1/3}$ 与 $x^{2/3}$ 的极值为什么不同？}
这里对负数也按实立方根理解：
\[
 r(x)=x^{1/3}=\sqrt[3]{x},\qquad q(x)=x^{2/3}=(\sqrt[3]{x})^2.
\]
\begin{center}% F002: cube root versus its square, exact same-scale plots.
\begin{tikzpicture}[font=\small,>=stealth,line cap=round]
\begin{scope}[shift={(0,0)}]
\node at (2.9,4.5) {$r(x)=\sqrt[3]{x}$};
\begin{scope}[shift={(2.9,2.2)},x=1.65cm,y=1.3cm]
  \draw[gray,->] (-1.35,0)--(1.4,0) node[right] {$x$};
  \draw[gray,->] (0,-1.3)--(0,1.2) node[below left=4pt] {$y$};
  \draw[blue!70!black,very thick,variable=\t,domain=-1.05:1.05,samples=121] plot ({\t*\t*\t},\t);
\draw[orange!85!black,thick] (0,-.42)--(0,.42);
  \fill[orange!85!black] (0,0) circle(2pt);
  \node[below right=5pt] at (0,0) {$O$};
\end{scope}
\node[align=center] at (2.9,-.6) {橙线为竖直切线：没有有限导数\\[5pt]仍一直增，$0$ 不是极值点};
\end{scope}
\begin{scope}[shift={(6.4,0)}]
\node at (2.9,4.5) {$q(x)=x^{2/3}$};
\begin{scope}[shift={(2.9,2.2)},x=1.65cm,y=1.3cm]
  \draw[gray,->] (-1.35,0)--(1.4,0) node[right] {$x$};
  \draw[gray,->] (0,-1.3)--(0,1.2) node[below left=4pt] {$y$};
  \draw[blue!70!black,very thick,variable=\t,domain=-1.05:1.05,samples=121] plot ({\t*\t*\t},{\t*\t});
\draw[orange!85!black,thick] (0,-.42)--(0,.42);
  \fill[orange!85!black] (0,0) circle(2pt);
  \node[below right=5pt] at (0,0) {$O$};
\end{scope}
\node[align=center] at (2.9,-.6) {尖点处没有有限导数\\[5pt]先降后升，$0$ 是极小值点};
\end{scope}
\end{tikzpicture}
\end{center}

\textbf{左图：不可导，没有极值。}当 $h\ne0$ 时，
\[
\frac{r(h)-r(0)}{h}=\frac{1}{(\sqrt[3]{h})^2}
\ \longrightarrow\ +\infty\qquad(h\to0).
\]
它在 $0$ 处没有有限导数，切线为竖直方向；但曲线仍从左向右一直增大。
特别地，原点左侧函数值小于 $0$、右侧大于 $0$，所以 $0$ 不是极值点。

\textbf{右图：不可导，却有极小值。}因为
\[
q(x)=(\sqrt[3]{x})^2>0=q(0)\qquad(x\ne0),
\]
所以 $0$ 是极小值点，极小值为 $0$，同时也是整个实数范围的最小值。
另一方面，
\[
\frac{q(h)-q(0)}{h}=\frac{1}{\sqrt[3]{h}}
\ \longrightarrow\
\begin{cases}-\infty,&h\to0^-,\\+\infty,&h\to0^+,\end{cases}
\]
也不能得到有限导数。左右两支在原点汇成尖点，曲线先降后升。

\textbf{结论：}两图都在原点不可导，但是否有极值仍由附近函数值决定；
$+\infty$ 或 $-\infty$ 都不能当作普通导数值。
前面四种情况中的 $|x|$，则是更简单的“不可导但有极小值”例子。
\end{explanationbox}
\clearpage
\begin{explanationbox}
\paragraph{极值只看附近，不要求在整个定义域最高或最低}

\begin{center}% F002: exact mathematical teaching diagram.
\begin{tikzpicture}[font=\small,>=stealth]
\begin{scope}[x=2.4cm,y=.65cm]
\fill[blue!7] (-1.35,-1.7) rectangle (-.65,6.6);
\draw[gray,->] (-2.2,0)--(2.5,0) node[right] {$x$};
\draw[gray,->] (0,-1.7)--(0,6.7) node[above] {$y$};
\draw[blue!70!black,very thick,domain=-2.1:2.3,samples=120] plot (\x,{\x*\x*\x-3*\x+1});
\draw[orange!85!black,very thick,domain=-1.35:-.65,samples=40] plot (\x,{\x*\x*\x-3*\x+1});
\draw[gray,dashed] (-1,0)--(-1,3)--(0,3);
\fill[orange!85!black] (-1,3) circle(2pt);
\node[above=10pt,align=center] at (-1,3) {附近最高\\$(-1,3)$};
\fill[blue!70!black] (2.3,6.267) circle(2pt);
\node[anchor=east,align=right] at (1.95,6.25) {远处可以更高\\$f(2.3)=6.267>3$};
\draw[gray,->] (1.95,6.25)--(2.22,6.267);
\node[below] at (-1,0) {$-1$};
\end{scope}
\node at (.5,5.1) {$f(x)=x^3-3x+1$：极大值只与附近比较};
\node at (.5,-1.85) {橙色一段及浅色带表示只看 $-1$ 附近的一小段。};
\end{tikzpicture}
\end{center}

图中 $x=-1$ 附近先增后减，$f(-1)=3$ 是极大值；右侧远处却有更高的函数值。
因此，\textbf{极大值是附近的最高值，不一定是整体最大值}；极小值也只和附近比较。

\paragraph{先固定比较范围：极值与最值用的尺子不同}
最大值、最小值统称\textbf{最值}；极大值、极小值统称\textbf{极值}。
设所讨论的比较范围为 $D$：极值只和判断点\textbf{附近}比较，最值要和\textbf{$D$ 中所有位置}比较，而且必须真正取到。

若 $D$ 上存在最大值 $M$、最小值 $m$，则
\[
 m\le f(x)\le M\qquad(\text{每一个 }x\in D).
\]
因此，\textbf{极大值不超过最大值，极小值不低于最小值}；它们可能相等，也可能不同。
上图的极大值 $3$ 就不是整个实数范围的最大值；实际上这条三次函数在整个实数范围内\textbf{既无最大值，也无最小值}。

\textbf{水平切线不等于转头。}上一节四图中，$x^3$ 和 $-x^3$ 在原点的切线也水平，
但一个一直增、一个一直减，都没有极值。
\end{explanationbox}
\clearpage
\begin{explanationbox}
\paragraph{内部的最值与极值可以重合，端点最值须另看}
\begin{center}% F002: same-scale quadratics on the same closed interval.
\begin{tikzpicture}[font=\small,>=stealth,line cap=round]
\node at (6.1,4.8) {同一比较范围 $[-1,1]$：看清内部与端点};
\node at (2.9,4.05) {$f(x)=-x^2$};
\node at (9.3,4.05) {$g(x)=x^2$};
\foreach \cx in {2.9,9.3}{
  \begin{scope}[shift={(\cx,1.85)},x=1.65cm,y=1.4cm]
    \draw[gray,->] (-1.3,0)--(1.4,0) node[right] {$x$};
    \draw[gray,->] (0,-1.35)--(0,1.2) node[below left=4pt] {$y$};
    \fill[orange!85!black] (0,0) circle(2pt);
    \node[below right=4pt] at (0,0) {$O$};
    \node[below=4pt] at (-1,0) {$-1$};
    \node[below=4pt] at (1,0) {$1$};
  \end{scope}
}
\begin{scope}[shift={(2.9,1.85)},x=1.65cm,y=1.4cm]
  \draw[blue!70!black,very thick,domain=-1:1,samples=81] plot (\x,{-\x*\x});
  \fill[blue!70!black] (-1,-1) circle(2pt);
  \fill[blue!70!black] (1,-1) circle(2pt);
\end{scope}
\begin{scope}[shift={(9.3,1.85)},x=1.65cm,y=1.4cm]
  \draw[blue!70!black,very thick,domain=-1:1,samples=81] plot (\x,{\x*\x});
  \fill[blue!70!black] (-1,1) circle(2pt);
  \fill[blue!70!black] (1,1) circle(2pt);
\end{scope}
\node[align=center] at (2.9,-.9)
  {内部：$0$ 既是极大值，也是最大值\\[5pt]端点：最小值 $-1$，不算内部极小值};
\node[align=center] at (9.3,-.9)
  {内部：$0$ 既是极小值，也是最小值\\[5pt]端点：最大值 $1$，不算内部极大值};
\end{tikzpicture}
\end{center}

图中两个内部的最值点，都比附近其他点\textbf{严格}更高或更低，所以同时是相应的极值点。
两个端点虽然取得最值，但本条只把\textbf{区间内部、左右都能比较}的位置称为极值点，因此不把端点算进去。

\textbf{不能省掉严格比较这个条件。}本文前面采用“附近其他点总比它小或大”的严格约定。
若最大值在内部取到，且附近其他点都\textbf{严格更小}，才同时是本文的极大值；最小值的说法相反。
\begin{center}% F002: a constant function attains global extrema, but no strict local extrema.
\begin{tikzpicture}[font=\small,>=stealth]
\node at (6,3.55) {$f(x)=1$，$-1\le x\le1$：最值存在，严格极值不存在};
\begin{scope}[shift={(6,.65)},x=3.6cm,y=1.75cm]
  \draw[gray,->] (-1.3,0)--(1.35,0) node[right] {$x$};
  \draw[gray,->] (0,-.22)--(0,1.25);
  \node[gray,left=6pt] at (0,.5) {$y$};
  \draw[blue!70!black,very thick] (-1,1)--(1,1);
  \fill[blue!70!black] (-1,1) circle(2pt);
  \fill[blue!70!black] (1,1) circle(2pt);
  \fill[orange!85!black] (0,1) circle(2pt);
  \node[above left=6pt] at (-1,1) {$1$};
  \node[below=5pt] at (-1,0) {$-1$};
  \node[below=5pt] at (1,0) {$1$};
\end{scope}
\node at (6,-.55) {最大值 $=$ 最小值 $=1$；任一点附近都有同高点，不能严格比较。};
\end{tikzpicture}
\end{center}
常数函数内部也处处取得最值，但附近有同高点，所以\textbf{不符合本文采用的严格极值定义}。
\end{explanationbox}
\clearpage
\begin{explanationbox}
\paragraph{极大值一定比极小值大吗？不同位置不能这样比}
\begin{center}% F002: exact trig model, radians converted explicitly for PGF trigonometry.
\begin{tikzpicture}[font=\small,>=stealth]
\node at (6,7.5) {$f(x)=x+2\sin x$：较早的极大值，可以小于较后的极小值};
\begin{scope}[shift={(.45,.65)},x=.88cm,y=.45cm]
  \draw[gray,->] (-.2,0)--(13.1,0) node[right] {$x$};
  \draw[gray,->] (0,-.6)--(0,13.6) node[below left=4pt] {$y$};
  \draw[blue!70!black,very thick,domain=0:12.5663706144,samples=201]
    plot (\x,{\x+2*sin(deg(\x))});
  \coordinate (A) at ({2*pi/3},{2*pi/3+sqrt(3)});
  \coordinate (B) at ({10*pi/3},{10*pi/3-sqrt(3)});
  \fill[orange!85!black] (A) circle(2.2pt);
  \fill[orange!85!black] (B) circle(2.2pt);
  \node[align=center] at ({2*pi/3},6.5) {极大值约 $3.826$\\$x=2\pi/3$};
  \draw[gray,->] ({2*pi/3},5.2)--({2*pi/3},4.05);
  \node[align=center] at ({10*pi/3},5.2) {极小值约 $8.740$\\$x=10\pi/3$};
  \draw[gray,->] ({10*pi/3},6.6)--({10*pi/3},8.5);
\end{scope}
\node at (6,-.55) {只与各自附近比较，不能把不同位置的极大值、极小值统一排序。};
\end{tikzpicture}
\end{center}

图中取 $f(x)=x+2\sin x$，有 $f'(x)=1+2\cos x$。在 $x=2\pi/3$ 处，导数由正变负，取得极大值
\[
 f(2\pi/3)=2\pi/3+\sqrt3\approx3.826;
\]
在 $x=10\pi/3$ 处，导数由负变正，取得极小值
\[
 f(10\pi/3)=10\pi/3-\sqrt3\approx8.740.
\]
这里\textbf{前一个极大值反而小于后一个极小值}。
这些数分别只和\textbf{各自附近}比较，不是同一范围里的整体最高、最低。
数值由上述精确式计算，四舍五入到小数点后三位。
\end{explanationbox}


\clearpage
\begin{explanationbox}
\paragraph{极值点说“在哪里”，极值说“多大”}

\begin{center}% F002: exact mathematical teaching diagram.
\begin{tikzpicture}[font=\small,>=stealth]
\node at (5.7,5.4) {$f(x)=-(x-2)^2+5$};
\begin{scope}[x=2.4cm,y=.72cm]
\draw[gray,->] (-.25,0)--(4.3,0) node[right] {$x$};
\draw[gray,->] (0,-.3)--(0,6.1) node[above] {$y$};
\draw[blue!70!black,very thick,domain=.5:3.5,samples=80] plot (\x,{-(\x-2)^2+5});
\draw[gray,dashed] (2,0)--(2,5)--(0,5);
\fill[orange!85!black] (2,5) circle(2pt);
\node[above=10pt] at (2,5) {$(2,5)$};
\node[below=5pt,align=center] at (2,0) {$2$\\极大值点：横坐标};
\node[left=8pt] at (0,5) {$5$};
\end{scope}
\node[anchor=west] at (9.1,4.0) {极大值为 $5$};
\node at (5.7,-1.9) {$x\ne2$ 时，$f(x)-f(2)=-(x-2)^2<0$};
\end{tikzpicture}
\end{center}

图中峰顶是曲线上的点 $(2,5)$：\textbf{极大值点的横坐标是 $2$，极大值是 $5$}。
虚线分别把位置与高度投影到两条坐标轴上。答题时不能把它们混写。
\end{explanationbox}
\clearpage
\begin{explanationbox}
\paragraph{一个实用的观察：减去函数值，再寻找平方因子}

\begin{center}% F002: exact mathematical teaching diagram.
\begin{tikzpicture}[font=\small,>=stealth]
\node at (5.8,4.0) {$q(x)=x^2(x+3)$：只看附近，就能判断极小值};
\begin{scope}[shift={(8.6,1.9)},x=2cm,y=.14cm]
\fill[teal!8] (-1,-18) rectangle (1,10);
\draw[gray,->] (-4.3,0)--(1.7,0) node[right] {$x$};
\draw[gray,->] (0,-18)--(0,10.7) node[below left=4pt] {$y$};
\draw[blue!70!black,very thick,domain=-4:1.35,samples=120] plot (\x,{\x*\x*(\x+3)});
\draw[orange!85!black,very thick,domain=-1:1,samples=60] plot (\x,{\x*\x*(\x+3)});
\fill[orange!85!black] (0,0) circle(2pt);
\node[below right=6pt] at (0,0) {$O(0,0)$};
\draw[gray,dashed] (-1,0)--(-1,2);
\draw[gray,dashed] (1,0)--(1,4);
\node[below=5pt] at (-1,0) {$-1$};
\node[below=5pt] at (1,0) {$1$};
\fill[blue!70!black] (-4,-16) circle(2pt);
\node[anchor=west,align=left] at (-2.8,-13) {远处可以更低\\$q(-4)=-16$};
\draw[gray,->] (-2.85,-13)--(-3.92,-15.8);
\end{scope}
\node at (5.8,-1.7) {$0<|x|<1$ 时，$x^2>0,\ x+3>2$，故 $q(x)>q(0)$。};
\end{tikzpicture}
\end{center}

先看具体例子 $q(x)=x^2(x+3)$。在 $0$ 附近，除了 $x=0$ 本身，
平方因子为正，$x+3$ 也为正，因此附近的曲线都在 $q(0)=0$ 上方。
远处即使出现更小的函数值，也不影响 $0$ 是极小值点。

一般地，如果
\[
f(x)-f(a)=(x-a)^2A(x),
\]
只要在 $a$ 附近 $A(x)>0$，就有 $f(x)>f(a)$（$x\ne a$），因此 $a$ 是极小值点；
若附近 $A(x)<0$，则是极大值点。\textbf{不必证明整个定义域内都更大或更小。}

\textbf{尖角也可能是最低处。}例如 $|x|$ 在 $0$ 处不可导，却在这里取得极小值；
后面的例4与易错一会直接画出这种情况。因此还要留意不能求导的内部点。
\end{explanationbox}
